\section{BLR over General Finite Fields}
\label{sec:blr-general-finite-fields}

Throughout this section, let $q=p^s$ for a prime $p$ and $s\in\bbN$ with $s\ge 1$.
All vectors are in $\bbF_q^n$, and all expectations are uniform over the indicated finite sets.
Let $\omega_p = e^{2\pi i/p}$.

\begin{definition}[Field trace]
\label{def:field-trace}
The field trace $\Tr: \bbF_q \to \bbF_p$ is
\[
    \Tr(z) := \sum_{i=0}^{s-1} z^{p^i}.
\]
\end{definition}

\begin{definition}[Additive characters]
\label{def:additive-character}
For $\alpha\in\bbF_q^n$, define the additive character
$\chi_\alpha: \bbF_q^n\to\bbC$ by
\[
    \chi_\alpha(x) := \omega_p^{\Tr(\inner{\alpha}{x})},
\]
where $\inner{\alpha}{x}=\sum_{i=1}^n \alpha_i x_i\in\bbF_q$.
\end{definition}

\begin{definition}[Inner product and Fourier coefficient]
\label{def:fourier-coefficient}
For functions $h_1,h_2:\bbF_q^n\to\bbC$, define
\[
    \inner{h_1}{h_2}
    := \bbE_{x\in\bbF_q^n}\left[h_1(x)\overline{h_2(x)}\right].
\]
For $h:\bbF_q^n\to\bbC$ and $\alpha\in\bbF_q^n$, define
\[
    \widehat h(\alpha)
    := \inner{h}{\chi_\alpha}
    = \bbE_{x\in\bbF_q^n}\left[h(x)\overline{\chi_\alpha(x)}\right].
\]
\end{definition}

\begin{lemma}[Character orthogonality]
\label{lem:character-orthogonality}
For every $\alpha,\beta\in\bbF_q^n$,
\[
    \bbE_{x\in\bbF_q^n}\left[\chi_\alpha(x)\overline{\chi_\beta(x)}\right]
    =
    \begin{cases}
        1 & \text{if } \alpha=\beta,\\
        0 & \text{if } \alpha\ne\beta.
    \end{cases}
\]
Equivalently, $\bbE_x[\chi_\alpha(x)]=0$ for $\alpha\ne 0$ and equals $1$ for $\alpha=0$.
\end{lemma}

\begin{proof}
First observe that for every $\alpha,\beta\in\bbF_q^n$ and every $x\in\bbF_q^n$,
\[
    \chi_\alpha(x)\overline{\chi_\beta(x)}
    =
    \omega_p^{\Tr(\inner{\alpha}{x})}
    \omega_p^{-\Tr(\inner{\beta}{x})}
    =
    \omega_p^{\Tr(\inner{\alpha-\beta}{x})}
    =
    \chi_{\alpha-\beta}(x),
\]
where we used the $\bbF_p$-linearity of the trace map.

If $\alpha=\beta$, then $\alpha-\beta=0$, so $\chi_{\alpha-\beta}(x)=1$ for every
$x\in\bbF_q^n$. Hence
\[
    \bbE_{x\in\bbF_q^n}\left[\chi_\alpha(x)\overline{\chi_\beta(x)}\right]
    =
    \bbE_{x\in\bbF_q^n}[1]
    =
    1.
\]

Now suppose $\alpha\ne\beta$. Let $\gamma:=\alpha-\beta$. Then $\gamma\ne 0$, so there exists
an index $j$ such that $\gamma_j\ne 0$. Fix all coordinates of $x$ except $x_j$. For the fixed
remaining coordinates, the map
\[
    x_j\mapsto \inner{\gamma}{x}
\]
has the form
\[
    x_j\mapsto \gamma_j x_j + c
\]
for some constant $c\in\bbF_q$. Since $\gamma_j\ne 0$, multiplication by $\gamma_j$ is a
bijection on $\bbF_q$, and therefore the affine map $x_j\mapsto \gamma_jx_j+c$ is also a
bijection on $\bbF_q$. Thus, as $x_j$ ranges uniformly over $\bbF_q$, the value
$\inner{\gamma}{x}$ also ranges uniformly over $\bbF_q$. Averaging first over $x_j$ and then
over the remaining coordinates gives
\[
    \bbE_{x\in\bbF_q^n}\left[\chi_\gamma(x)\right]
    =
    \bbE_{t\in\bbF_q}\left[\omega_p^{\Tr(t)}\right].
\]

It remains to show that this last average is zero. The trace map
$\Tr:\bbF_q\to\bbF_p$ is a nonzero $\bbF_p$-linear map, so every fiber of $\Tr$ has the same
cardinality. Since $\abs{\bbF_q}=q$ and $\abs{\bbF_p}=p$, each fiber has cardinality $q/p$.
Therefore
\begin{align*}
    \bbE_{t\in\bbF_q}\left[\omega_p^{\Tr(t)}\right]
    &=
    \frac{1}{q}\sum_{t\in\bbF_q}\omega_p^{\Tr(t)} \\
    &=
    \frac{1}{q}\sum_{r\in\bbF_p}
        \sum_{\substack{t\in\bbF_q\\ \Tr(t)=r}}
        \omega_p^r \\
    &=
    \frac{1}{q}\sum_{r\in\bbF_p}\frac{q}{p}\omega_p^r \\
    &=
    \frac{1}{p}\sum_{r\in\bbF_p}\omega_p^r \\
    &=0,
\end{align*}
because the sum of all $p$-th roots of unity is zero. Hence, if $\alpha\ne\beta$, then
\[
    \bbE_{x\in\bbF_q^n}\left[\chi_\alpha(x)\overline{\chi_\beta(x)}\right]
    =
    0.
\]

Combining the two cases proves the stated orthogonality relation. The equivalent formulation
follows by taking $\beta=0$.
\end{proof}

\begin{lemma}[Fourier inversion and Parseval]
\label{lem:fourier-inversion-parseval}
For every function $h:\bbF_q^n\to\bbC$,
\[
    h(x)=\sum_{\alpha\in\bbF_q^n}\widehat h(\alpha)\chi_\alpha(x)
    \quad\text{for every }x\in\bbF_q^n,
\]
and
\[
    \sum_{\alpha\in\bbF_q^n}\abs{\widehat h(\alpha)}^2
    = \bbE_{x\in\bbF_q^n}\left[\abs{h(x)}^2\right].
\]
In particular, if $\abs{h(x)}=1$ for every $x$, then
$\sum_{\alpha\in\bbF_q^n}\abs{\widehat h(\alpha)}^2=1$.
\end{lemma}

\begin{proof}
We use the character basis of the additive group \(F^{\mathrm{Idx}}\). For
each \(\alpha \in F^{\mathrm{Idx}}\), let \(\chi_\alpha\) denote the additive
character
\[
    \chi_\alpha(x)
    =
    \omega_p^{\operatorname{Tr}(\langle \alpha,x\rangle)} .
\]
By the orthogonality of additive characters,
\[
    \langle \chi_\alpha,\chi_\beta\rangle
    =
    \mathbb E_x\!\left[\chi_\alpha(x)\overline{\chi_\beta(x)}\right]
    =
    \begin{cases}
    1, & \alpha=\beta,\\
    0, & \alpha\neq \beta.
    \end{cases}
\]
Moreover, these characters span the vector space of all complex-valued
functions on \(F^{\mathrm{Idx}}\). Hence
\(\{\chi_\alpha : \alpha \in F^{\mathrm{Idx}}\}\) is an orthonormal basis with
respect to the inner product
\[
    \langle g,h\rangle
    =
    \mathbb E_x\!\left[g(x)\overline{h(x)}\right].
\]

Therefore every \(g : F^{\mathrm{Idx}}\to \mathbb C\) admits a unique expansion
\[
    g
    =
    \sum_{\alpha \in F^{\mathrm{Idx}}} c_\alpha \chi_\alpha .
\]
Taking the inner product with \(\chi_\beta\) gives
\[
    \langle g,\chi_\beta\rangle
    =
    \sum_{\alpha \in F^{\mathrm{Idx}}}
        c_\alpha \langle \chi_\alpha,\chi_\beta\rangle
    =
    c_\beta .
\]
By definition, \(\widehat g(\beta)=\langle g,\chi_\beta\rangle\). Hence
\(c_\beta=\widehat g(\beta)\) for every \(\beta\), and so
\[
    g
    =
    \sum_{\alpha \in F^{\mathrm{Idx}}}
        \widehat g(\alpha)\chi_\alpha .
\]
Evaluating this identity at \(x\) gives the Fourier inversion formula.

For Parseval's identity, use the same expansion:
\[
\begin{aligned}
    \|g\|_2^2
    &= \langle g,g\rangle \\
    &=
    \left\langle
        \sum_{\alpha}\widehat g(\alpha)\chi_\alpha,
        \sum_{\beta}\widehat g(\beta)\chi_\beta
    \right\rangle \\
    &=
    \sum_{\alpha,\beta}
        \widehat g(\alpha)
        \overline{\widehat g(\beta)}
        \langle \chi_\alpha,\chi_\beta\rangle .
\end{aligned}
\]
By orthonormality, all terms with \(\alpha\neq\beta\) vanish, while
\(\langle\chi_\alpha,\chi_\alpha\rangle=1\). Therefore
\[
    \|g\|_2^2
    =
    \sum_{\alpha}
        \widehat g(\alpha)\overline{\widehat g(\alpha)}
    =
    \sum_{\alpha} |\widehat g(\alpha)|^2 .
\]
Since \(\|g\|_2^2=\mathbb E_x |g(x)|^2\), this proves Parseval's identity.
\end{proof}

\begin{definition}[Fourier expansion of a finite-field-valued function]
\label{def:fourier-expansion-field-valued}
Let $f:\bbF_q^n\to\bbF_q$. For each $c\in\bbF_q^\times$, define
\[
    \fexp{\ff}{c}(x) := \omega_p^{\Tr(c f(x))}.
\]
The Fourier expansion of $f$ is the family
\[
    \left\{\fexp{\ff}{c}\right\}_{c\in\bbF_q^\times}.
\]
\end{definition}

\begin{lemma}[Character-sum indicator]
\label{lem:character-sum-indicator}
For every $z\in\bbF_q$,
\[
    \ind[z=0]
    = \frac{1}{q}\sum_{c\in\bbF_q}\omega_p^{\Tr(cz)}.
\]
Consequently, for every $u,v\in\bbF_q$,
\[
    \ind[u=v]
    = \frac{1}{q}\sum_{c\in\bbF_q}\omega_p^{\Tr(c(u-v))}.
\]
\end{lemma}
\begin{proof}
Let
\[
    \psi(z) := \omega_p^{\Tr(z)}
\]
be the standard additive character of \(\bbF_q\). We prove that, for every
\(t\in \bbF_q\),
\[
    \ind[t=0]
    =
    \frac{1}{q}\sum_{c\in\bbF_q}\psi(ct).
\]
The desired statement follows by taking \(t=u-v\).

If \(t=0\), then \(ct=0\) for every \(c\in\bbF_q\), so
\[
    \frac{1}{q}\sum_{c\in\bbF_q}\psi(ct)
    =
    \frac{1}{q}\sum_{c\in\bbF_q}1
    =
    1.
\]

Now suppose \(t\neq 0\). Then multiplication by \(t\) is a bijection
\(\bbF_q\to\bbF_q\), so
\[
    \sum_{c\in\bbF_q}\psi(ct)
    =
    \sum_{z\in\bbF_q}\psi(z).
\]
We claim that the latter sum is \(0\). Let
\[
    S:=\sum_{z\in\bbF_q}\psi(z).
\]
Since \(\Tr:\bbF_q\to\bbF_p\) is a nonzero \(\bbF_p\)-linear map, it is
surjective. Hence there exists \(a\in\bbF_q\) such that \(\Tr(a)=1\).
Using the bijection \(z\mapsto z+a\), we get
\[
    S
    =
    \sum_{z\in\bbF_q}\psi(z+a)
    =
    \sum_{z\in\bbF_q}\omega_p^{\Tr(z+a)}
    =
    \sum_{z\in\bbF_q}\omega_p^{\Tr(z)}\omega_p^{\Tr(a)}
    =
    \omega_p S.
\]
Since \(\omega_p\neq 1\), this implies \(S=0\). Therefore, when
\(t\neq 0\),
\[
    \frac{1}{q}\sum_{c\in\bbF_q}\psi(ct)=0.
\]

Combining the two cases gives
\[
    \ind[t=0]
    =
    \frac{1}{q}\sum_{c\in\bbF_q}\omega_p^{\Tr(ct)}.
\]
Substituting \(t=u-v\), we obtain
\[
    \ind[u=v]
    =
    \frac{1}{q}\sum_{c\in\bbF_q}\omega_p^{\Tr(c(u-v))}.
\]
\end{proof}

\begin{definition}[Relative Hamming distance]
\label{def:relative-hamming-distance}
For $f,g:\bbF_q^n\to\bbF_q$, define
\[
    \delta(f,g)
    := \Pr_{x\in\bbF_q^n}\left[f(x)\ne g(x)\right].
\]
\end{definition}

\begin{definition}[Linear functions]
\label{def:linear-functions}
For $\alpha\in\bbF_q^n$, define $\ell_\alpha:\bbF_q^n\to\bbF_q$ by
\[
    \ell_\alpha(x):=\inner{\alpha}{x}.
\]
The set of linear functions is
\[
    \lin := \left\{\ell_\alpha \mid \alpha\in\bbF_q^n\right\}.
\]
The distance from $f:\bbF_q^n\to\bbF_q$ to linearity is
\[
    \delta(f,\lin):=\min_{\alpha\in\bbF_q^n}\delta(f,\ell_\alpha).
\]
\end{definition}

\begin{lemma}[Distance formula]
\label{lem:distance-formula}
Let $f,g:\bbF_q^n\to\bbF_q$.
Let $\{\fexp{\ff}{c}\}_{c\in\bbF_q^\times}$ and
$\{\fexp{\fg}{c}\}_{c\in\bbF_q^\times}$ be their Fourier expansions.
Then
\[
    \delta(f,g)
    = 1-\frac{1}{q}\left(1+
        \sum_{c\in\bbF_q^\times}
        \inner{\fexp{\ff}{c}}{\fexp{\fg}{c}}
    \right).
\]
Equivalently,
\[
    \delta(f,g)
    = 1-\frac{1}{q}\left(1+
        \sum_{c\in\bbF_q^\times}
        \sum_{\alpha\in\bbF_q^n}
        \widehat{\fexp{\ff}{c}}(\alpha)
        \overline{\widehat{\fexp{\fg}{c}}(\alpha)}
    \right).
\]
\end{lemma}

\begin{proof}
By \cref{lem:character-sum-indicator},
\[
    \Pr_x[f(x)=g(x)]
    = \bbE_x\left[\frac{1}{q}\sum_{c\in\bbF_q}
        \omega_p^{\Tr(c(f(x)-g(x)))}\right].
\]
The term $c=0$ contributes $1/q$. For $c\ne 0$,
\[
    \omega_p^{\Tr(c(f(x)-g(x)))}
    = \fexp{\ff}{c}(x)\overline{\fexp{\fg}{c}(x)}.
\]
Therefore
\[
    \Pr_x[f(x)=g(x)]
    =\frac{1}{q}\left(1+\sum_{c\in\bbF_q^\times}
        \inner{\fexp{\ff}{c}}{\fexp{\fg}{c}}
    \right),
\]
which gives the first formula after using
$\delta(f,g)=1-\Pr_x[f(x)=g(x)]$.
The second formula follows from Fourier inversion and Parseval.
\end{proof}

\begin{lemma}[Fourier expansion of a linear function]
\label{lem:linear-function-fourier-expansion}
Fix $\alpha\in\bbF_q^n$ and $c\in\bbF_q^\times$.
Let $\fexp{\lambda}{c}(x):=\omega_p^{\Tr(c\ell_\alpha(x))}$.
Then
\[
    \fexp{\lambda}{c}=\chi_{c\alpha}.
\]
Equivalently, for every $\beta\in\bbF_q^n$,
\[
    \widehat{\fexp{\lambda}{c}}(\beta)
    =
    \begin{cases}
        1 & \text{if } \beta=c\alpha,\\
        0 & \text{if } \beta\ne c\alpha.
    \end{cases}
\]
\end{lemma}

\begin{proof}
For every $x\in\bbF_q^n$,
\[
    \fexp{\lambda}{c}(x)
    = \omega_p^{\Tr(c\inner{\alpha}{x})}
    = \omega_p^{\Tr(\inner{c\alpha}{x})}
    = \chi_{c\alpha}(x).
\]
The statement about Fourier coefficients follows from
\cref{lem:character-orthogonality}.
\end{proof}

\begin{lemma}[Distance to linearity in Fourier form]
\label{lem:distance-to-linearity-fourier}
Let $f:\bbF_q^n\to\bbF_q$ have Fourier expansion
$\{\fexp{\ff}{c}\}_{c\in\bbF_q^\times}$. Then
\[
    \delta(f,\lin)
    = 1-\frac{1}{q}\left(1+
        \max_{\alpha\in\bbF_q^n}
        \sum_{c\in\bbF_q^\times}
        \widehat{\fexp{\ff}{c}}(c\alpha)
    \right).
\]
Moreover, for every $\alpha\in\bbF_q^n$, the quantity
\[
    S_\alpha := \sum_{c\in\bbF_q^\times}
        \widehat{\fexp{\ff}{c}}(c\alpha)
\]
is real and satisfies $-1\le S_\alpha\le q-1$.
\end{lemma}

\begin{proof}
Apply \cref{lem:distance-formula} with $g=\ell_\alpha$ and use
\cref{lem:linear-function-fourier-expansion}. This gives
\[
    \delta(f,\ell_\alpha)
    =1-\frac{1}{q}\left(1+
        \sum_{c\in\bbF_q^\times}
        \widehat{\fexp{\ff}{c}}(c\alpha)
    \right).
\]
Taking the minimum over $\alpha$ gives the claimed formula for
$\delta(f,\lin)$.
For fixed $\alpha$, the same formula expresses $S_\alpha$ as
$q(1-\delta(f,\ell_\alpha))-1$. Hence $S_\alpha$ is real, and since
$0\le \delta(f,\ell_\alpha)\le 1$, we get $-1\le S_\alpha\le q-1$.
\end{proof}

\begin{definition}[Finite-field BLR test]
\label{def:finite-field-blr-test}
Given oracle access to $f:\bbF_q^n\to\bbF_q$, the finite-field BLR verifier
$\BLRVerifier$ samples $a,b\in\bbF_q^\times$ and $x,y\in\bbF_q^n$ uniformly and accepts iff
\[
    a f(x)+b f(y)=f(ax+by).
\]
That is,
\[
    \BLRVerifier=1
    \quad\Longleftrightarrow\quad
    a f(x)+b f(y)=f(ax+by).
\]
\end{definition}

\begin{lemma}[Completeness]
\label{lem:blr-completeness}
If $f\in\lin$, then
\[
    \Pr[\BLRVerifier=1]=1.
\]
\end{lemma}

\begin{proof}
If $f=\ell_\alpha$ for some $\alpha\in\bbF_q^n$, then for every
$a,b\in\bbF_q^\times$ and $x,y\in\bbF_q^n$,
\[
    f(ax+by)
    =\inner{\alpha}{ax+by}
    =a\inner{\alpha}{x}+b\inner{\alpha}{y}
    =a f(x)+b f(y).
\]
Thus the verifier accepts for every choice of randomness.
\end{proof}

\begin{lemma}[Exact BLR acceptance formula]
\label{lem:exact-blr-acceptance}
Let $f:\bbF_q^n\to\bbF_q$ have Fourier expansion
$\{\fexp{\ff}{c}\}_{c\in\bbF_q^\times}$.
For each $\alpha\in\bbF_q^n$, define
\[
    S_\alpha := \sum_{c\in\bbF_q^\times}
        \widehat{\fexp{\ff}{c}}(c\alpha).
\]
Then
\[
    \Pr[\BLRVerifier=1]
    = \frac{1}{q}\left(1+\frac{1}{(q-1)^2}
        \sum_{\alpha\in\bbF_q^n} S_\alpha^3
    \right).
\]
\end{lemma}

\begin{proof}
By \cref{lem:character-sum-indicator}, for fixed $a,b,x,y$,
\[
    \ind[a f(x)+b f(y)=f(ax+by)]
    =\frac{1}{q}\sum_{c\in\bbF_q}
        \omega_p^{\Tr(c(a f(x)+b f(y)-f(ax+by)))}.
\]
Taking expectation gives
\[
    \Pr[\BLRVerifier=1]
    =\frac{1}{q}+\frac{1}{q}\bbE_{a,b,x,y}
        \left[\sum_{c\in\bbF_q^\times}
        \fexp{\ff}{ca}(x)\fexp{\ff}{cb}(y)
        \fexp{\ff}{-c}(ax+by)
        \right].
\]
Expand the three complex-valued functions into their Fourier expansions:
\begin{align*}
&\bbE_{a,b,x,y}
        \left[\sum_{c\in\bbF_q^\times}
        \fexp{\ff}{ca}(x)\fexp{\ff}{cb}(y)
        \fexp{\ff}{-c}(ax+by)
        \right] \\
&=\bbE_{a,b}\left[\sum_{c\in\bbF_q^\times}
    \sum_{\alpha,\beta,\gamma\in\bbF_q^n}
    \widehat{\fexp{\ff}{ca}}(\alpha)
    \widehat{\fexp{\ff}{cb}}(\beta)
    \widehat{\fexp{\ff}{-c}}(\gamma)
    \bbE_{x,y}\left[
        \chi_\alpha(x)\chi_\beta(y)\chi_\gamma(ax+by)
    \right]
\right].
\end{align*}
Since
\[
    \chi_\gamma(ax+by)=\chi_{a\gamma}(x)\chi_{b\gamma}(y),
\]
\cref{lem:character-orthogonality} implies that the inner expectation is $1$ exactly when
\[
    \alpha+a\gamma=0
    \quad\text{and}\quad
    \beta+b\gamma=0,
\]
and is $0$ otherwise. Therefore the last display equals
\[
    \bbE_{a,b}\left[\sum_{c\in\bbF_q^\times}
    \sum_{\alpha\in\bbF_q^n}
    \widehat{\fexp{\ff}{ca}}(\alpha)
    \widehat{\fexp{\ff}{cb}}(a^{-1}b\alpha)
    \widehat{\fexp{\ff}{-c}}(-a^{-1}\alpha)
    \right].
\]
Substitute $\alpha=ca\eta$. Then this becomes
\[
    \bbE_{a,b}\left[\sum_{c\in\bbF_q^\times}
    \sum_{\eta\in\bbF_q^n}
    \widehat{\fexp{\ff}{ca}}(ca\eta)
    \widehat{\fexp{\ff}{cb}}(cb\eta)
    \widehat{\fexp{\ff}{-c}}(-c\eta)
    \right].
\]
As $a,b,c$ range over $\bbF_q^\times$, so do $ca$, $cb$, and $-c$. Hence this equals
\[
    \frac{1}{(q-1)^2}
    \sum_{\eta\in\bbF_q^n}
    \left(\sum_{d\in\bbF_q^\times}
        \widehat{\fexp{\ff}{d}}(d\eta)
    \right)^3
    =\frac{1}{(q-1)^2}\sum_{\eta\in\bbF_q^n}S_\eta^3.
\]
Substituting into the expression for $\Pr[\BLRVerifier=1]$ proves the claim.
\end{proof}

\begin{lemma}[Second moment bound]
\label{lem:second-moment-bound}
With $S_\alpha$ as in \cref{lem:exact-blr-acceptance},
\[
    \sum_{\alpha\in\bbF_q^n} S_\alpha^2 \le (q-1)^2.
\]
\end{lemma}

\begin{proof}
Expanding the square gives
\[
    \sum_{\alpha\in\bbF_q^n} S_\alpha^2
    =\sum_{a,b\in\bbF_q^\times}
        \sum_{\alpha\in\bbF_q^n}
        \widehat{\fexp{\ff}{a}}(a\alpha)
        \widehat{\fexp{\ff}{b}}(b\alpha).
\]
For fixed $a,b\in\bbF_q^\times$, Cauchy--Schwarz and
\cref{lem:fourier-inversion-parseval} give
\begin{align*}
\abs{\sum_{\alpha\in\bbF_q^n}
        \widehat{\fexp{\ff}{a}}(a\alpha)
        \widehat{\fexp{\ff}{b}}(b\alpha)}
&\le
\left(\sum_{\alpha\in\bbF_q^n}\abs{\widehat{\fexp{\ff}{a}}(a\alpha)}^2\right)^{1/2}
\left(\sum_{\alpha\in\bbF_q^n}\abs{\widehat{\fexp{\ff}{b}}(b\alpha)}^2\right)^{1/2} \\
&=1.
\end{align*}
The equality uses that multiplication by $a$ and by $b$ permutes $\bbF_q^n$, and that
$\abs{\fexp{\ff}{a}(x)}=\abs{\fexp{\ff}{b}(x)}=1$ for every $x$.
Summing over the $(q-1)^2$ choices of $(a,b)$ gives the result.
\end{proof}

\begin{theorem}[Soundness of the finite-field BLR test]
\label{thm:finite-field-blr-soundness}
For every function $f:\bbF_q^n\to\bbF_q$,
\[
    \Pr[\BLRVerifier=0] \ge \delta(f,\lin).
\]
Equivalently,
\[
    \Pr[\BLRVerifier=1] \le 1-\delta(f,\lin).
\]
\end{theorem}

\begin{proof}
Let $S_\alpha$ be as in \cref{lem:exact-blr-acceptance}, and let
\[
    M:=\max_{\alpha\in\bbF_q^n}S_\alpha.
\]
By \cref{lem:distance-to-linearity-fourier},
\[
    1-\delta(f,\lin)=\frac{1}{q}(1+M).
\]
By \cref{lem:distance-to-linearity-fourier}, each $S_\alpha$ is real and $S_\alpha\le M$.
Therefore $S_\alpha^3\le M S_\alpha^2$ for every $\alpha$, because $S_\alpha^2\ge 0$.
Using \cref{lem:exact-blr-acceptance} and \cref{lem:second-moment-bound},
\begin{align*}
    \Pr[\BLRVerifier=1]
    &=\frac{1}{q}\left(1+\frac{1}{(q-1)^2}
        \sum_{\alpha\in\bbF_q^n}S_\alpha^3\right) \\
    &\le \frac{1}{q}\left(1+\frac{M}{(q-1)^2}
        \sum_{\alpha\in\bbF_q^n}S_\alpha^2\right) \\
    &\le \frac{1}{q}(1+M) \\
    &=1-\delta(f,\lin).
\end{align*}
Taking complements gives $\Pr[\BLRVerifier=0]\ge\delta(f,\lin)$.
\end{proof}
